\documentclass{article}
\usepackage[T1]{fontenc}
\usepackage[utf8]{inputenc}
\usepackage{amsmath}
\usepackage{comment}
\usepackage{amssymb}
\usepackage{polski}
\usepackage[polish]{babel}
\usepackage{graphicx}
\usepackage{float} % Inf
\usepackage{enumitem}
\usepackage{subcaption}
\usepackage{yfonts}
\usepackage{braket} % kwant
\usepackage{tikz}
\usetikzlibrary{positioning, arrows.meta, calc}
\usepackage[margin=2cm]{geometry}
\usepackage{tabularx}
%\usepackage[version=4]{mhchem} % Chemia
\usepackage{hyperref}
\hypersetup{
colorlinks=true,
linkcolor=black,
urlcolor=red
}
\usepackage{listings}
\lstset{
  basicstyle=\ttfamily\small,
  breaklines=false,
  frame=single
}

\title{\Huge{\textgoth{Sprawozdanie - Projekt 6}}}
\author{\Huge{\textfrak{Maciej Flaga}}}
\date{}

\begin{document}
\maketitle

\section{Realizacja}

Napisano funkcję \texttt{diagYslice}, która zwracała wektor wartości własnych Hamiltonianu układu jednowymiarowego będącego przekrojem głównego układu dla ustalonego \texttt{x[i]}. W tym celu skorzystano z biblioteki \texttt{GSL} i jej funkcji \texttt{gsl\_eigen\_symm} oraz \texttt{gsl\_sort\_vector}. Po zdiagonalizowaniu hamiltonianu dla każdego przekroju użyto funkcji napisanych tydzień temu, które rozwiązują układ metodą QTBM. Techniczny przebieg kodu:
\begin{lstlisting}
maciej@ThinkPadT490:~/kod/sem6/uklnw/lab6$ make
gcc main.c -Wall -Wextra -Werror -Wpedantic -O3 -lm -lgsl -lgslcblas -o main.out
./main.out
Wszystkie (nx=720) diagonalizacje w czasie 1.3539s
Wyznaczenie G metoda QTBM w czasie 1.4036s
python3 plot.py
Czas wykreslania: 0.9514s
\end{lstlisting}

\subsection{Wykresy}

\begin{figure}[H]
    \centering
    \includegraphics[width=1.15\linewidth]{Vmap.png}
    \caption{Mapa potencjału emulującego bramkę}
    % \label{fig:placeholder}
\end{figure}

\begin{figure}[H]
    \centering
    \includegraphics[width=1\linewidth]{Vnmap.png}
    \caption{Mapa energii własnych dla wszystkich przekrojów}
    % \label{fig:placeholder}
\end{figure}

\begin{figure}[H]
    \centering
    \includegraphics[width=0.72\linewidth]{Gfile.png}
    \caption{Wykres konduktancji od energii}
    % \label{fig:placeholder}
\end{figure}

\subsection{Wnioski}

Widzimy, że konduktancja układa się we wzór schodków. Zdecydowano się wykreślić znormalizowaną konduktancje $G/G_0$ tak aby pokazać kwanty konduktancji w całkowitych wartościach. Schodki są lekko zaokrąglone co może być spowodowane zastosowaniem przybliżenia adiabatycznego.

\end{document}
